\textbf{拉格朗日函数 / 对偶 / KKT 速查}

\textbf{术语}
\begin{itemize}
	\item 下凸 = convex， 上凸 = concave。
\end{itemize}

\textbf{标准形式}
\[
\begin{aligned}
\min_{\mathbf{x}\in\mathbb R^n}\quad & f(\mathbf{x}) \\
\text{s.t.}\quad & h_i(\mathbf{x})=0,\quad i=1,\dots,m \\
& g_j(\mathbf{x})\le 0,\quad j=1,\dots,p
\end{aligned}
\]

\textbf{拉格朗日函数}
\[
\mathcal L(\mathbf{x},\boldsymbol{\lambda},\boldsymbol{\mu})
=f(\mathbf{x})+\sum_{i=1}^m \lambda_i h_i(\mathbf{x})+\sum_{j=1}^p \mu_j g_j(\mathbf{x}),
\qquad \mu_j\ge 0
\]
\begin{itemize}
	\item $\lambda_i$ 无符号限制。
	\item 若原式写成 $g_j(\mathbf{x})\ge 0$，先改成 $-g_j(\mathbf{x})\le 0$。
\end{itemize}

\textbf{只有等式约束}
\[
\mathcal L(\mathbf{x},\boldsymbol{\lambda})=f(\mathbf{x})+\sum_{i=1}^m \lambda_i h_i(\mathbf{x})
\]
\[
\nabla_{\mathbf{x}}\mathcal L(\mathbf{x},\boldsymbol{\lambda})=0,\qquad h_i(\mathbf{x})=0
\]

\textbf{对偶函数}
\[
q(\boldsymbol{\lambda},\boldsymbol{\mu})
=\inf_{\mathbf{x}\in\mathcal D}\mathcal L(\mathbf{x},\boldsymbol{\lambda},\boldsymbol{\mu})
\]
\begin{itemize}
	\item 对任意 $\boldsymbol{\mu}\ge \mathbf{0}$，有 $q(\boldsymbol{\lambda},\boldsymbol{\mu})\le p^*$。
	\item $q$ 总是上凸函数。
\end{itemize}

\textbf{对偶问题}
\[
\max_{\boldsymbol{\lambda},\boldsymbol{\mu}}\ q(\boldsymbol{\lambda},\boldsymbol{\mu})
\qquad \text{s.t.}\quad \boldsymbol{\mu}\ge \mathbf{0}
\]
\[
d^*:=\max q(\boldsymbol{\lambda},\boldsymbol{\mu}),\qquad p^*:=\text{原问题最优值}
\]

\textbf{弱对偶 / 强对偶}
\[
d^*\le p^*
\]
\[
d^*=p^* \iff \text{强对偶}
\]

\textbf{Slater 条件}

原问题为下凸优化，且
\begin{itemize}
	\item $f,g_j$ 为下凸函数；
	\item $h_i$ 为仿射函数；
	\item 存在严格可行点 $\bar{\mathbf{x}}$：
	\[
	h_i(\bar{\mathbf{x}})=0,\qquad g_j(\bar{\mathbf{x}})<0
	\]
\end{itemize}

\textbf{常用结论}
\begin{itemize}
	\item 下凸问题 + Slater $\Rightarrow$ 强对偶。
	\item 常见有限维情形下，可取到最优乘子。
\end{itemize}

\textbf{鞍点}
\[
\mathcal L(\mathbf{x}^*,\boldsymbol{\lambda},\boldsymbol{\mu})
\le \mathcal L(\mathbf{x}^*,\boldsymbol{\lambda}^*,\boldsymbol{\mu}^*)
\le \mathcal L(\mathbf{x},\boldsymbol{\lambda}^*,\boldsymbol{\mu}^*),
\qquad \forall \mathbf{x},\boldsymbol{\lambda},\boldsymbol{\mu}\ (\boldsymbol{\mu}\ge \mathbf{0})
\]
\begin{itemize}
	\item 固定 $(\boldsymbol{\lambda}^*,\boldsymbol{\mu}^*)$，$\mathbf{x}^*$ 使 $\mathcal L$ 最小。
	\item 固定 $\mathbf{x}^*$，$(\boldsymbol{\lambda}^*,\boldsymbol{\mu}^*)$ 使 $\mathcal L$ 最大。
	\item 在下凸问题中，若满足 Slater 且最优解可取到，则
	\[
	\text{强对偶} \iff \mathcal L \text{ 存在鞍点}.
	\]
\end{itemize}

\textbf{KKT 条件}

若 $\mathbf{x}^*$ 为局部最优解，且满足约束资格条件，则存在 $\boldsymbol{\lambda}^*,\boldsymbol{\mu}^*$ 使
\begin{itemize}
	\item 平稳性
	\[
	\nabla f(\mathbf{x}^*)+\sum_{i=1}^m \lambda_i^* \nabla h_i(\mathbf{x}^*)+\sum_{j=1}^p \mu_j^* \nabla g_j(\mathbf{x}^*)=\mathbf{0}
	\]
	\item 原始可行性
	\[
	h_i(\mathbf{x}^*)=0,\qquad g_j(\mathbf{x}^*)\le 0
	\]
	\item 对偶可行性
	\[
	\boldsymbol{\mu}^*\ge \mathbf{0}
	\]
	\item 互补松弛
	\[
	\mu_j^*g_j(\mathbf{x}^*)=0
	\]
\end{itemize}

\textbf{互补松弛速记}
\begin{itemize}
	\item $g_j(\mathbf{x}^*)<0 \Rightarrow \mu_j^*=0$
	\item $\mu_j^*>0 \Rightarrow g_j(\mathbf{x}^*)=0$
\end{itemize}

\textbf{KKT 何时可直接判最优}
\begin{itemize}
	\item 一般非下凸问题：KKT 只是必要条件。
	\item 可微下凸优化 + Slater：KKT 也是充分条件。
	\item 此时 KKT / 强对偶 / 鞍点可视作同一组结论。
\end{itemize}

\textbf{最后速记}
\begin{itemize}
	\item 最小化 + $g\le 0$ $\Rightarrow \mu\ge 0$。
	\item 先写 $\mathcal L$，再写 KKT。
	\item 对偶是“最大下界”。
	\item Slater 用来保证强对偶。
	\item 非下凸时，KKT 不保证最优。
\end{itemize}
